\documentclass[openany]{book}

\usepackage[margin=1in]{geometry}
\usepackage{amsmath,amsfonts,amsthm, amssymb}
\usepackage{yhmath}
\usepackage{mathrsfs}
\usepackage{mathtools}
\usepackage{xcolor}
\usepackage{graphicx}
\usepackage{comment}
\usepackage{tikz-cd}
\usepackage{quiver}
\usepackage{hyperref,cleveref}
\renewcommand{\familydefault}{ppl}
\newcommand{\curl}{\text{curl}}
\newcommand{\diverg}{\text{div}}
\newcommand{\tr}{\text{tr}}
\newcommand{\R}{\mathbb{R}}
\newcommand{\E}{\mathbb{E}}
\newcommand{\Z}{\mathbb{Z}}
\newcommand{\C}{\mathbb{C}}
\newcommand{\F}{\mathbb{F}}
\newcommand{\la}{\langle}
\newcommand{\ra}{\rangle}
\newcommand{\colim}{\text{colim}}
\DeclareMathOperator{\im}{im}
\DeclareMathOperator{\disc}{disc}
\let\oldemptyset\emptyset
\let\emptyset\varnothing
\newcommand{\tor}{\text{Tor}}
\newcommand{\id}{\text{id}}
\newcommand{\ext}{\text{Ext}}
\newcommand{\ptop}{\text{PTop}}
\newcommand{\pt}{\text{pt}}
\newcommand{\ach}{\text{Ach}}
\newcommand{\Q}{\mathbb{Q}}
\newcommand{\gal}{\text{Gal}}
\newcommand{\fraccomma}{\genfrac{}{}{0pt}{}{}{,}}
\definecolor{wikipediadarkblue}{rgb}{0.023, 0.270, 0.676}
\hypersetup{
    colorlinks,
    citecolor=black,
    filecolor=black,
    linkcolor=blue,
    urlcolor=red
}

% \hypersetup{
%     urlbordercolor={1 0 0}, % Red box for URLs
%     linkbordercolor=blue, % Blue box for internal links
%     % the values are in RGB format from 0 to 1
% }


\input{hui_r.tex}

\title{Complex Analysis Qualifying Exam Review
\\ 
\vspace{0.4cm}
\large Spring 2026}



% \author{(Please email hsun95@jh.edu if you see typos)}
% \date{\today}


\begin{document}

\maketitle

\tableofcontents
\newpage



\chapter{Basic calculus in the complex domain}

\begin{prop}[Cauchy-Riemann Equations]
    If $f:\Omega\to\C$ is holomorphic, then $\frac{\partial f}{\partial x}, \frac{\partial f}{\partial y}$ exist and are continuous on $\Omega$, and 
    \begin{equation*}
        \frac{\partial f}{\partial x}=\frac{1}{i}\frac{\partial f}{\partial y}
    \end{equation*}
    If we write $f=u+iv$, then the above is equivalent to 
    \begin{equation*}
        \begin{cases}
            \frac{\partial u}{\partial x}=\frac{\partial v}{\partial y}\\
            \frac{\partial v}{\partial x}=-\frac{\partial u}{\partial y}
        \end{cases}
    \end{equation*}
\end{prop}

\begin{prop}[Converse of Cauchy-Riemann]
    If $f:\Omega\to\C$ is $C^1$ and satisfies Cauchy-Riemann equations, then $f$ is holomorphic.
\end{prop}

% \begin{prop}
%     If $f:\overline{\Omega}\to\C$ is holomorphic, where $z_0\in\Omega$ and $D_r(z_0)\subset\Omega$, then for $z\in D_r(z_0)$, $f$ is given by a convergent power series 
%     \begin{equation*}
%         f(z)=\sum_{k=0}^\infty a_k(z-z_0)^k
%     \end{equation*}
%     where 
%     \begin{equation*}
%         a_k=\frac{1}{2\pi i}\int_{\partial\Omega}\frac{f(\zeta)}{(\zeta-z_0)^{k+1}}d\zeta=\frac{f^{(k)}(z_0)}{k!}
%     \end{equation*}
% \end{prop}
% (This follows from Cauchy integral formula).


% \begin{defn}[complex log]
%     The logorithmic function extended to the complex plane is $\log:\C\setminus\R^-\to\C$, given by 
%     \begin{equation*}
%         \log z=\frac{1}{z}\frac{1}{\zeta}d\zeta
%     \end{equation*}
% \end{defn}


\begin{prop}
    Let $f:\Omega\to\C$ be holomorphic, where $\Omega$ is open. Let $p\in\Omega$, assume $f'(p)\neq 0$. Then there exists a neighborhood $O$ of $p$ and a neighborhood $U$ of $q=f(p)$ such that $f:O\to U$ is bijective, where the inverse $g=f^{-1}: U\to O$ is holomorphic. For $z\in O, w=f(z)$, we have 
    \begin{equation*}
        g'(w)=\frac{1}{f'(z)}
    \end{equation*}
\end{prop}


\begin{prop}
    Suppose $\Omega\subset\C$ is convex. Assume $f$ is holomorphic in $\Omega$ and there exists $a\in\C$ such that 
    \begin{equation*}
        \text{Re } af'>0 \text{ on } \Omega
    \end{equation*}
    then $f$ is one-to-one onto its image $f(\Omega)$.
\end{prop}




\chapter{Cauchy integral theorem}

\begin{thm}
    If $f\in C^1(\overline{\Omega})$ is holomorphic on $\Omega$, then 
    \begin{equation*}
        \int_{\partial\Omega}f(z)dz=0
    \end{equation*}
\end{thm}

\begin{thm}
    If $f\in C^1(\overline{\Omega})$ is holomorphic on $\Omega$ and $z_0\in\Omega$, then 
    \begin{equation*}
        f(z_0)=\frac{1}{2\pi i}\int_{\partial\Omega}\frac{f(z)}{z-z_0}dz
    \end{equation*}
    One can change $\partial\Omega$ to a simple closed curve $\gamma$ completely contained in $\Omega$.
\end{thm}


\begin{thm}[holomorphic implies analytic]
    If $f\in C^1(\overline{\Omega})$ is holomorphic, then for $z\in D_r(z_0)\subset\Omega$, then $f(z)$ has the convergent power series expansion 
    \begin{equation*}
        f(z)=\sum_{n=0}^\infty a_n(z-z_0)^n
    \end{equation*}
    where 
    \begin{equation*}
        a_n=\frac{1}{2\pi i}\int_{\partial\Omega}\frac{f(\zeta)}{(\zeta-z_0)^{n+1}}d\zeta=\frac{f^{(n)}(z_0)}{n!}
    \end{equation*}
\end{thm}

\begin{prop}
    If $f$ is holomorphic on an open set $\Omega\subset\C$, then 
    \begin{equation*}
        f\in C^\infty(\Omega)
    \end{equation*}
\end{prop}



\begin{prop}
    Let $\Omega\subset\C$ be an open set and let $f_\nu:\Omega\to\C$ be holomorphic. Assume $f_\nu\to f$ locally uniformly, then $f:\Omega\to\C$ is holomorphic (and $f_\nu'\to f'$ locally uniformly on $\Omega$).
\end{prop} 


\begin{prop}[maximum principle]
    Let $\Omega\subset\C$ be connected, open. If $f$ is holomorphic in $\Omega$, then for $z_0\in\Omega$,
    \begin{equation*}
        \text{ if } |f(z_0)|=\sup_{z\in\Omega}|f(z)|, \text{ then } f \text{ is constant on $\Omega$}
    \end{equation*}
    Moreover, if $\Omega$ is bounded and $f\in C(\overline{\Omega})$, then 
    \begin{equation*}
        \sup_{z\in\overline{\Omega}}|f(z)|=\sup_{z\in\partial\Omega}|f(z)|
    \end{equation*}
    (In this case, $f$ must attain a maximum in $\overline{\Omega}$). This basically states that either $f$ is constant or $f$ attains maximum on the boundary.
\end{prop}


% \begin{prop}
%     Suppose $f$ is holomorphic on the unit disk $D_1(0)$. Assume $|f(z)|\leq 1$ for $|z|<1$, and $f(0)=0$. Then 
%     \begin{equation*}
%         |f(z)|\leq |z|
%     \end{equation*}
%     Moreover, $|f(z)|=|z|$ iff $f(z)=cz$ for some $c$ such that $|c|=1$.
% \end{prop}

\begin{thm}[Liouville's theorem]
    If $f:\C\to\C$ is holomorphic and bounded, then $f$ is constant.
\end{thm}
A consequence of Liouville's theorem is the fundamental theorem of algebra:
\begin{prop}
    If $p(z)$ is a polynomial of degree $n\geq 1$, then $p(z)$ has at least one root. (Proof: consider $1/p(z)$).
\end{prop}
There are several properties that also hold for harmonic functions.
\begin{prop}[mean value property of harmonic functions]
    If $u\in C^2(\Omega)$ is harmonic, $z_0\in\Omega$, and $\overline{D_r(z_0)}\subset\Omega$, then 
    \begin{equation*}
        u(z_0)=\frac{1}{2\pi}\int_0^{2\pi}u(z_0+re^{i\theta})d\theta
    \end{equation*}
\end{prop}


\begin{prop}
    Let $\Omega\subset\C$ be connected and open. If $u:\Omega\to\R$ is harmonic in $\Omega$, then for $z_0\in\Omega$,
    \begin{equation*}
        u(z_0)=\sup_{z\in\Omega}u(z)\Rightarrow u \text{ is constnat on $\Omega$}
    \end{equation*}
    Moreover, if $\Omega$ is bounded, $u\in C(\overline{\Omega})$, then 
    \begin{equation*}
        \sup_{z\in\overline{\Omega}}u(z)=\sup_{z\in\partial\Omega}u(z)
    \end{equation*}
\end{prop}

\begin{prop}[Liouville for harmonic functions]
    If $u\in C^2(\C)$ is bounded and harmonic on all of $\C$, then $u$ is constant.
\end{prop}

\begin{prop}
    Let $\Omega\subset\C$ be simply connected domain. Then each harmonic function $u\in C^2(\Omega)$ has a harmonic conjugate.
\end{prop}
Next is a converse of Cauchy's integral theorem.
\begin{thm}[Morera's theorem]
    Let $g:\Omega\to\C$ be continuous, and 
    \begin{equation*}
        \int_\gamma g(z)dz=0
    \end{equation*}
    for any $\gamma=\partial R$, where $R\subset\Omega$ is a rectangle. Then $g$ is holomorphic.
\end{thm}


\begin{prop}[Schwarz reflection principle]
    Assume $f:\Omega^+\cup L\to\C$ is continuous, holomorphic in $\Omega^+$, and real-valued on $L$. Define $g:\Omega\to\C$ by
    \begin{equation*}
        g(z)=f(z), z\in\Omega^+\cup L, \overline{f(\overline{z})}, z\in\Omega^-
    \end{equation*}
    Then $g$ is holomorphic on $\Omega$.
\end{prop}


\begin{prop}[Goursat's theorem]
    If $f:\Omega\to\C$ is complex-differentiable at each point of $\Omega$, then $f$ is holomorphic.
\end{prop}



\begin{prop}[analytic continuation]
    Let $\Omega \subset \mathbb{C}$ be open and connected, and let $f : \Omega \to \mathbb{C}$ be holomorphic. If $f = 0$ on a set $S \subset \Omega$ and $S$ has an accumulation point $p \in \Omega$, then $f = 0$ on $\Omega$.

    In other words, if $f=g$ on a set $S\subset\Omega$, where $S$ contains an accumulation point, then $f=g$ on $\Omega$.
\end{prop}



% \begin{prop}
%     Let $\Omega \subset \mathbb{C}$ be a simply connected domain. Assume $f : \Omega \to \mathbb{C}$ is holomorphic and nowhere vanishing. Then there exists a holomorphic function $g$ on $\Omega$ such that
%     \begin{equation*}
%         e^{g(z)} = f(z), \quad \forall z \in \Omega.
%         \end{equation*}
% \end{prop}


\begin{defn}[singularities]
    There are three types of singularities: $p$ is said to be the following for $f:U\to\C$,
    \begin{enumerate}
        \item Removable: if there exists $\tilde{f}$ on $U$ such that $\tilde{f}=f$ on $U\setminus p$.
        \item Pole: if $|f(z)|\to\infty$ as $z\to p$. Equivalently, $p$ is a removable singularity for $1/f$.
        \item Essential singularity: not removable, not pole, an isolated singularity. $f(z)=e^{1/z}$.
    \end{enumerate}
\end{defn}


\begin{prop}
    If $p\in\Omega$ and $f$ is holomorphic on $\Omega\setminus p$ and bounded, then $p$ is a removable singularity.
\end{prop}

\begin{prop}
    If $f$ is holomorphic on $\Omega\setminus p$ with a pole at $p$, then there exists $k\in\mathbb{Z}^+$ such that 
    \begin{equation*}
        f(z)=\frac{F(z)}{(z-p)^k}
    \end{equation*}
    on $\Omega\setminus p$, where $F$ is holomorphic on $\Omega$ and $F(p)\neq 0$.
\end{prop}


\begin{prop}[Casorati-Weierstrass]
    Suppose $f:\Omega\setminus p\to\C$ has an essential singularity at $p$. Then for any neighborhood of $U$ of $p$ in $\Omega$, the image of $U\setminus p$ is dense in $\C$.
\end{prop}


\begin{prop}
    If $f:\C\to\C$ is holomorphic and $|f(z)|\to\infty$ as $|z|\to\infty$, then $f(z)$ is a polynomial.
\end{prop}

\begin{prop}
    Let $f:\mathcal{A}\to\C$ be holomorphic, where the annulus $\mathcal{A}=\{z: r_0\leq |z-z_0|\leq r_1\}$, then $f$ is given by 
    \begin{equation*}
        f(z)=\sum_{n=-\infty}^\infty a_n(z-z_0)^n, z\in\mathcal{A} \tag{\text{Laurent series}}
    \end{equation*}
    and 
    \begin{equation*}
        a_n=\frac{1}{2\pi i}\int_\gamma \frac{f(\zeta)}{(\zeta-z_0)^{n+1}}d\zeta, n\in\mathbb{Z}
    \end{equation*}
    where $\gamma$ is any counterclockwise circle centered at $z_0$, of radius $r\in(r_0,r_1)$. 
\end{prop}
How does the Laurent series expansion help us?
\begin{prop}
    Suppose $f$ is holomorphic on $D_R(z_0)\setminus z_0$, with the above Laurent series. $f$ has a pole at $z_0$ iff $a_n=0$ for all but finitely many $n<0$; $f$ has an essential singularity iff $a_n\neq 0$ for infinitely many $n$.
\end{prop}



\chapter{Fourier analysis and complex function theory}


\begin{prop}
    Let $f\in C(S^1)$, if the Fourier coefficients $\hat{f}(k)$ form a summable series, i.e., if 
    \begin{equation*}
        \sum_{k=-\infty}^\infty|\hat{f}(k)|<\infty
    \end{equation*}
    then for each $\theta\in S^1$, 
    \begin{equation*}
        f(\theta)=\sum_{k=-\infty}^\infty\hat{f}(k)e^{ik\theta}
    \end{equation*}
\end{prop}



\begin{prop}
    If $f$ is such that $\sum|\hat{f}(k)|<\infty$, then 
    \begin{equation*}
        \sum|\hat{f}(k)|^2=\frac{1}{2\pi}\int_{S^1}|f(\theta)|^2d\theta
    \end{equation*}
    More generally, if also $g$ is such that $\sum|\hat{g}(k)|<\infty$, 
    \begin{equation*}
        \sum\hat{f}(k)\overline{\hat{g}(k)}=\frac{1}{2\pi}\int_{S^1}f(\theta)\overline{g(\theta)}d\theta
    \end{equation*}
\end{prop}

\begin{prop}
    If $f$ is a continuous, piecewise $C^1$ function on $S^1$, then $\sum|\hat{f}(k)|<\infty$.
\end{prop}

\begin{prop}
    If $f$ is Riemann integrable on $S^1$, then $S_Nf\to f$ in $L^2$, and 
    \begin{equation*}
        \sum|\hat{f}(k)|^2=\|f\|_{L^2}^2
    \end{equation*}
\end{prop}

\begin{prop}[Dirichlet boundary problem]
    If $D\subset\C$ is an open disk and given $f\in C(\partial D)$, there exists a unique $u\in C(\overline{D})\cap C^2(D)$ satisfying 
    \begin{equation*}
        \Delta u=0 \text{ on } D, u\vert_{\partial D}=f
    \end{equation*}
\end{prop}


\begin{prop}
    Assume $u:\Omega^+\cup L\to\C$ is continuous, harmonic on $\Omega^+$, and $u=0$ on $L$. Define $v:\Omega\to\C$ as 
    \begin{equation*}
        \begin{cases}
            v(z)=u(z), z\in \Omega^+\cup L\\
            -u(\bar{z}), z\in\Omega^-
        \end{cases}
    \end{equation*}
    then $v$ is harmonic on $\Omega$.
\end{prop}



\begin{prop}
    If $f$ is integrable on $\R$, then $\hat{f}$ is continuous on $\R$.
\end{prop}
The next proposition gives criterion for the Fourier inversion formula.
\begin{prop}
    Assume $f$ is bounded and continuous on $\R$, and $f,\hat{f}$ are integrable on $\R$. Then the Fourier inversion formula holds: for all $x\in\R$,
    \begin{equation*}
        f(x)=\frac{1}{\sqrt{2\pi}}\int_{-\infty}^\infty\hat{f}(\xi)e^{ix\xi}d\xi
    \end{equation*}
\end{prop}
We define 
\begin{equation*}
    S_Rf(x)=\frac{1}{\sqrt{2\pi}}\int_{-R}^R\hat{f}(\xi)e^{ix\xi}d\xi
\end{equation*}
\begin{prop}
    If $f\in L^2(\R)$, then as $R\to\infty$,
    \begin{equation*}
        S_Rf\to f \text{ in } L^2
    \end{equation*}
\end{prop}


\begin{prop}[uniqueness]
    If $f\in L^1(\R)$ and $\hat{f}=0$, then 
    \begin{equation*}
        f=0
    \end{equation*}
\end{prop}


\begin{defn}
    Let $f:\R^+\to\C$ that is integrable on $[0,R]$, for all $R<\infty$, such that 
    \begin{equation*}
        \int_0^\infty|f(t)|e^{-at}dt<\infty
    \end{equation*}
    for all $a>A$, for some $A\in(-\infty,\infty)$. Then we define the Laplace transform of $f$ 
    \begin{equation*}
        \mathcal{L}f(s)=\int_0^\infty f(t)e^{-st}dt, \text{ Re} s>A
    \end{equation*}
\end{defn}





\chapter{Residue calculus, the argument principle, and two very special
functions}

\begin{defn}
    If $f$ has the Laurent series in disk $D_j$ around $p_j\in\Omega$,
    \begin{equation*}
        f(z)=\sum_{n=-\infty}^\infty a_n(p_j)(z-p_j)^n
    \end{equation*}
    where the coefficient $a_{-1}(p_j)$ is called the residue of $f$ at $p_j$, denoted as $\text{Res}_{p_j}f$. We have 
    \begin{equation*}
        \text{Res}_{p_j}f=\frac{1}{2\pi i}\int_{\partial D_j}f(z)dz
    \end{equation*} 

\end{defn}


\begin{prop}[Residue formula]
    Let $f$ be a meromorphic function on $\Omega$, then for a simple closed curve $C$, we have 
    \begin{equation*}
        \int_Cf(z)dz=2\pi i\sum_{j=1}^k\text{Res}(f,z_j)
    \end{equation*}
    where the isolated singularities $z_j$'s are contained in $C$.
\end{prop}


\begin{prop}[Residue]
    If $f$ has a simple pole at $z_0$, then 
    \begin{equation*}
        \text{Res}(f,z_0)=\lim_{z\to z_0}(z-z_0)f(z)
    \end{equation*}
    If $f$ has a pole of order $n$ at $z_0$, then
    \begin{equation*}
        \text{Res}(f,z_0)=\lim_{z\to z_0}\frac{d^{n-1}}{dz^{n-1}}\left((z-z_0)^nf(z)\right)
    \end{equation*}
    If $f$ has an essential singularity at $z_0$, then we write $f$ as Laurent series around $z_0$,
    \begin{equation*}
        f(z)=\sum_{n=-\infty}^\infty c_n(z-z_0)^n
    \end{equation*}
    where 
    \begin{equation*}
        \text{Res}(f,z_0)=c_{-1}
    \end{equation*}
\end{prop}

\begin{prop}
    Let $f\in C^2(\bar{\Omega})$ be holomorphic in $\Omega$, nowhere zero on $\partial\Omega$, then let $\nu(f,\Omega)$ be the number of zeros of $f$ in $\Omega$, counted with multiplicity, 
    \begin{equation*}
        \nu(f,\Omega)=\frac{1}{2\pi i}\int_{\partial\Omega}\frac{f'(z)}{f(z)}dz
    \end{equation*}
\end{prop}



\begin{prop}
    Let $C_j$ denote the connected component of $\partial\Omega$, and $\gamma_j=f(C_j)$. Then the winding number of $\gamma_j$ around $0$ satisfies
    \begin{equation*}
        \frac{1}{2\pi i}\int_{\gamma_j}\frac{1}{z}dz=\frac{1}{2\pi i}\int_{C_j}\frac{f'(z)}{f(z)}dz
    \end{equation*}
\end{prop}

\begin{defn}[winding number]
    Let $\gamma$ be a curve given by 
    \begin{equation*}
        \gamma(t)=r(t)e^{i\theta(t)}
    \end{equation*}
    where $r(t),\theta(t)$ both continuous, piecewise $C^1$, and $r(t)>0$. The \textbf{winding number} of $\gamma$ about $0$ is given by 
    \begin{equation*}
        n(\gamma,0)=\frac{1}{2\pi i}\int_\gamma\frac{1}{z}dz
    \end{equation*}
\end{defn}

\begin{prop}
    If $\gamma_0,\gamma_1$ are smoothly homotopic in $\C\setminus 0$, then 
    \begin{equation*}
        n(\gamma_0,0)=n(\gamma_1,0)
    \end{equation*}
\end{prop}
Combining all the above,
\begin{equation*}
    \nu(f,\Omega)=\sum_j n(\gamma_j, 0)
\end{equation*}
the number of zeros of $f$, counted with multiplicity, is equal to the sum of the winding numbers of $f(C_j)$ about $0$.
\begin{prop}[Argument principle for meromorphic functions]
    If $f$ is meromorphic on a bounded domain $\Omega$ and $C^1$ on a neighborhood of $\partial\Omega$. Then the number of zeros of $f$ minus the number of poles, counted with multiplicity is equal to the sum of winding numbers of $f(C_j)$ around 0, where $C_j$ are the connected components of $\partial\Omega$.
\end{prop}

\begin{thm}[Rouche's theorem]
    Let $f,g\in C^1(\bar{\Omega})$ be holomorphic in $\Omega$, and nowhere zero on $\partial\Omega$. Assume 
    \begin{equation*}
        |f(z)-g(z)|<|f(z)|, \text{ for all } z\in\partial\Omega
    \end{equation*}
    then 
    \begin{equation*}
        \nu(f,\Omega)=\nu(g,\Omega)
    \end{equation*}
\end{thm}


\begin{prop}[Open mapping theorem]
    If $\Omega\subset\C$ is open and connected and $f:\Omega\to\C$ is holomorphic and non-constant, then $f$ maps open sets to open sets.
\end{prop}


\begin{defn}[Gamma function]
    The gamma function $\Gamma$ is given by 
    \begin{equation*}
        \Gamma(z)=\int_0^\infty e^{-t}t^{z-1}dt, \text{Re}z>0
    \end{equation*}
\end{defn}

\begin{defn}[Riemann zeta function]
    The zeta function is 
    \begin{equation*}
        \zeta(s)=\sum_{n=1}^\infty n^{-s}, \text{Re}s>1
    \end{equation*}
\end{defn}


\begin{prop}
    For $z\in\C\setminus\Z$, we have 
    \begin{equation*}
        \Gamma(z)\Gamma(1-z)=\frac{\pi}{\sin\pi z}
    \end{equation*}
\end{prop}
\begin{prop}
    $\Gamma(z)$ has no zeros, so $1/\Gamma(z)$ is an entire function.
\end{prop}

\begin{prop}
    For $\text{Re}z>0$, we have 
    \begin{equation*}
        \Gamma(z)=\lim_{n\to\infty}n^z\frac{1\cdot 2\cdot\dots \cdot n}{z(z+1)\dots(z+n)}
    \end{equation*}
\end{prop}


\begin{prop}[Euler's product formula]
    Let $\{p_j:j\geq 1\}$ denote the set of prime numbers, then for $\text{Re}s>1$, 
    \begin{equation*}
        \zeta(s)=\prod_{j=1}^\infty\frac{1}{1-p_j^{-s}}
    \end{equation*}
\end{prop}

\begin{prop}
    The function $\zeta(s)$ extends to a meromorphic function on $\C$, with one simple pole at $s=1$. Moreover, 
    \begin{equation*}
        \zeta(s)-\frac{1}{s-1}
    \end{equation*}
    extends to an entire function of $s$.
\end{prop}






\chapter{Conformal maps and geometrical aspects of complex function theory}

\begin{defn}[conformal]
    We say $f$ is conformal at $z_0$ if $\gamma_1, \gamma_2$ are two smooth paths that intersect at $z_0$, then the images $f(\gamma_j)$ intersect at $f(z_0)$ at the same angle as $\gamma_1$ and $\gamma_2$.
\end{defn}


\begin{prop}
    A smooth function $f$ on $\Omega$ is holomorphic iff it is conformal and preserves orientation.
\end{prop}



\begin{thm}[Riemann mapping theorem]
    If $\Omega\subset\C$ is a simply connected domain and $\Omega\neq\C$, then there exists a holomorphic diffeomorphism
    \begin{equation*}
        F:\Omega\to D
    \end{equation*}
    where $D$ is the unit disk. (A holomorphic diffeomorphism is a holomorphic function that has a holomorphic inverse).
\end{thm}
The following is an improvement of Casorati-Weierstrass theorem.
\begin{thm}[Picard's theorem]
    Let $p,q$ be distinct and 
    \begin{equation*}
        f:D\setminus\{0\}\to\C\setminus\{p,q\}
    \end{equation*}    
    is holomorphic, then the singularity at $0$ is either a pole or a removable singularity.
\end{thm}


\begin{prop}[Schwarz Lemma]
    Let $f:\mathbb{D}\to\mathbb{D}$ is holomorphic with $f(0)=0$, then 
    \begin{enumerate}
        \item $|f(z)|\leq |z|$ for all $z\in\mathbb{D}$.
        \item For some $z_0\neq 0$, if $|f(z_0)|=|z_0|$, then $f$ is a rotation.
        \item We have $|f'(0)|\leq 1$, and if it attains equality, then it implies $f$ is a rotation.
    \end{enumerate}
\end{prop}







% \chapter{Useful facts}
% \begin{enumerate}
%     \item $e^z$ has the following power series expansion: 
%     \begin{equation*}
%         e^z=\sum_{j=0}^\infty\frac{z^j}{j!}
%     \end{equation*}
%     \item $\cos, \sin$ has the following power series expansions:
%     \begin{equation*}
%         \cos z=
%     \end{equation*}
%     \begin{equation*}
%         \sin z=
%     \end{equation*}
% \end{enumerate}







% \chapter{Maybe}
% \begin{prop}
%     Let $\Omega$ be simply connected. If $1\in\Omega, 0\not\in\Omega$, then it has a branch of $\log$, where $F(z)=\log z$ such that 
%     \begin{enumerate}
%         \item $F$ is holomorphic.
%         \item $e^{F(z)}=z$ for all $z\in\Omega$.
%     \end{enumerate}
%     On $\C\setminus(-\infty,0]$, where 
%     \begin{equation*}
%         \log z=\log r+i\theta
%     \end{equation*}
%     is called the principal branch.
% \end{prop}



\end{document}